Heterogena robotflottor i lager består av agenter med kompletterande roller, såsom transportfordon och plockrobotar, vars beslut måste samordnas med hänsyn till uppgiftsberoenden, gemensamma resurser och batteribegränsningar. Detta examensarbete undersöker om lokalt körda, allmänt tillgängliga \glspl{llm} kan stödja beslutsfattande på hög nivå för denna form av samarbete mellan flera robotar. En grundad styrloop implementerades i simuleringsmiljön Battery-TA-RWARE: lagrets aktuella tillstånd omvandlas till en prompt, modellen väljer diskreta makrohandlingar och de föreslagna handlingarna valideras mot miljöns genomförbarhetsvillkor före exekvering.

Experimenten omfattar sex språkmodeller med ungefär 1 till 14 miljarder parametrar och sex lagerscenarier. De jämför centraliserad gemensam planering med delad kontext och separat planering för varje agent, representation i naturligt språk med \gls{json}-baserad representation samt tre promptbaserade målfamiljer som behandlar hyllleveranser, användning av \gls{llm}-anrop och batterihantering. Prestandan bedöms med hjälp av bland annat antal hyllleveranser, anropsfrekvens, ett leveransmått normaliserat efter antal modellanrop och laddningsbeteende. Fel vid generering eller handlingsvalidering loggas för diagnostisk validering av körningarna.

Resultaten visar att planeringsarkitektur, promptkonfiguration, modellval och scenariokonfiguration påverkade samarbetets prestanda. Konfigurationen med naturligt språk och planeringsarkitekturen med delad kontext var generellt förknippade med bättre leveransresultat i den utvärderade miljön, även om den senare krävde fler modellanrop. Bland de utvalda modellartefakterna ökade prestandan inte konsekvent med den nominella parameterstorleken, och förändringar av det angivna målet gav olika beteenden hos olika modeller.

Resultaten indikerar att språkmodeller kan fatta grundade beslut inom ett begränsat styrgränssnitt för heterogena robotar, men också att ett effektivt samarbete beror på samspelet mellan modellen och det omgivande planeringssystemet. Eftersom antalet körningar och initialiseringar är begränsat tolkas resultaten som deskriptiva och motiverar bredare utvärderingar av språkmodellbaserat robotsamarbete.
